% A REGISTERED RESULT SECTION. No numeral outside \ccnum; see check_numbers.py.
\section{What building it taught}
\label{sec:lessons}

Each paragraph is one lesson from building the tool: what happened, and what I changed. The most costly one, checks that passed without running, is in Section~\ref{sec:blind-spots}.

\paragraph{The models moved faster than I could measure.} After a change to the merge prompt, every merge I sent to \texttt{\ccnum{opus5.id}} through its API was refused. The vendor's name for the reason was \texttt{reasoning\_extraction}. Going back through the versions of the prompt, I found the one paragraph that caused it, a note that the sources are shown one segment per line. Without it the request was answered, and other models of the same vendor answered it with it. I dropped the model on \ccnum{opus5.retired_on} for its successor, \texttt{\ccnum{opus5.replaced_by}}. The requests of that search are published in \texttt{\ccnum{reg.2026-09-25-vendor.path}}.

The program around a model moved too. After an update of the vendor's command-line program, one model stopped looking facts up at the level \texttt{sourced}, and only the report, not the exit code, showed it. A \texttt{sourced} merge that retrieved nothing now ends as inconclusive. My answer to both was to keep the tool and the measurement apart: the tool has versions and is tested on its own, and a benchmark is made with one fixed set of models and carries its date.

\paragraph{A nested command inherits its user's configuration.} The subscription route runs Claude Code as a program inside the tool. Started plainly, it loaded my own configuration, including my instruction file \texttt{CLAUDE.md}, which is written for coding sessions and could shape a merge. The route now starts the program in its safe mode with only the tools the task needs, and the report records both. The benchmark of Section~\ref{sec:findings} ran in safe mode throughout. My earlier subscription figures were made with my configuration loaded and should not be compared with it.

\paragraph{Answer keys come from a diff, not from reading.} The Voyager set was first scored against a list of errors that was typed by reading each source beside its reference. A word-by-word comparison of the same files found many more. Reading had found the loud errors and missed the quiet ones, a swapped word or a decimal comma. Those are exactly the errors a model carries through unnoticed. The answer keys are now built by the published scoring program and were reviewed error by error.

\paragraph{A claim cannot be graded by the words it is made of.} I wanted a program, with no model, to decide questions of meaning: whether a merge really merged its sources or pasted them together, whether a replacement carries what the missing sentence said, whether a title the merge wrote agrees with the sources. Each attempt was a measure of the text alone with a cut-off: similarity, the proportion of lengths, the count of missing segments, the words a title shares with the sources. On the test documents the values for correct and for defective merges overlapped, so no cut-off separates them:
\begin{points}
\item \emph{A merge that says everything twice.} I compared the sentences of a merge with each other. The most similar pair of sentences in a document that is known to be correct was more similar than any pair in the merges that really stated both sources' facts twice. A cut-off that catches those merges accuses the correct document first.
\item \emph{A merge that only pastes.} Similarity to the pasted sources fired on a nearly complete copy. It missed a merge that had pasted both sources and dropped their titles and summary lines.
\item \emph{The length of the merge.} The proportion of lengths was higher for a merge with no repeated fact than for one with several, so it does not measure repetition.
\item \emph{A replacement that does not carry its segment.} Similarity between a missing segment and its replacement flags \ccnum{b5.reference_percent}\% of the correct records in the reference merges, and counting every missing segment against the loss budget puts a correct reference merge over the budget (Section~\ref{sec:declared-blindspot}).
\item \emph{A title.} The rule was that every substantial word of a rewritten title must occur in the sources. The example in the test suite: a source says ``a quad feels more stable at walking pace'', and the title is ``A motorbike is more stable than a quad''. Every word of that title is in the sources, so it passes the rule, and it turns the statement round.
\end{points}

\noindent A measure of the surface answers a different question, and its answer is worse than none, because it reads like a verification. So these questions are left to the judge. The fact stated twice, for example, is now found among the claims that the judge draws from the merge. The one rule of this kind that survived does not read at all. It counts records (Section~\ref{sec:declared-blindspot}), and it narrows its hole without closing it.
